% Lapisan solusi O001 -- Bab 9 (Ruang Vektor Topologis)
% 1 solusi asli terpisah; bukan tulisan John M. Erdman.
% Provenans model: OpenAI Codex gpt-5.6-sol, Ultra, atas arahan pengguna.
% Lisensi komponen solusi: CC BY-SA 4.0.
% Tidak menyiratkan dukungan John M. Erdman atau Portland State University.

\markboth{SOLUSI O001 UNTUK BAB 9: RUANG VEKTOR TOPOLOGIS}{SOLUSI O001 UNTUK BAB 9: RUANG VEKTOR TOPOLOGIS}
\section*{Solusi O001 untuk Bab 9: Ruang Vektor Topologis}
\addcontentsline{toc}{section}{Solusi O001 untuk Bab 9: Ruang Vektor Topologis}
\begin{center}
\small
Komponen solusi asli terpisah, CC BY-SA 4.0. Ditulis dengan bantuan
OpenAI Codex gpt-5.6-sol, Ultra, atas arahan pengguna. Pernyataan latihan
berasal dari edisi terjemahan karya John M. Erdman; solusi ini bukan tulisan
beliau dan tidak menyiratkan dukungan beliau atau Portland State University.
\end{center}

% O001-SOLUTION-ID: O001-FAOA-2015-CH09-EX-001-SOLUTION
% SOURCE-EXERCISE-ID: FAOA-2015-CH09-NODE-0072
% STATEMENT-TARGET-SHA256: 69f9a89b1b499eb3fee7d7d92605c324b0dabd8f66918754428f27cd5482fdfe
\begin{o001solution}
  {O001-FAOA-2015-CH09-EX-001-SOLUTION}
  {FAOA-2015-CH09-NODE-0072}
  {69f9a89b1b499eb3fee7d7d92605c324b0dabd8f66918754428f27cd5482fdfe}
\begin{o001statement}
Jelaskan pula bagaimana proposisi sebelumnya memungkinkan kita mengaitkan suatu ruang lain yang bersifat Hausdorff dengan
suatu ruang vektor topologis non-Hausdorff.
\end{o001statement}
\begin{o001answer}
Untuk ruang vektor topologis $X$, ruang yang terkait secara kanonik adalah
Hausdorffisasi
\[
  X_{\mathrm H}:=X/\overline{\{\vc0\}}.
\]
Ruang ini Hausdorff, dan setiap pemetaan linear kontinu dari $X$ ke ruang
vektor topologis Hausdorff berfaktor secara unik melalui pemetaan hasil bagi
$X\to X_{\mathrm H}$.
\end{o001answer}
\begin{o001proof}
Misalkan $N=\overline{\{\vc0\}}$. Karena penjumlahan, invers aditif, dan
perkalian skalar kontinu, tutupan subruang vektor tetap merupakan subruang
vektor; jadi $N$ adalah subruang tertutup dari~$X$. Proposisi sebelumnya
langsung memberi bahwa ruang hasil bagi $X/N$ dengan topologi hasil bagi
bersifat Hausdorff. Hal ini tetap berlaku ketika $X$ sendiri non-Hausdorff;
dalam kasus tersebut $N\ne\{\vc0\}$.

Konstruksi ini kanonik. Jika $Y$ ruang vektor topologis Hausdorff dan
$T:X\to Y$ linear kontinu, maka $\{\vc0\}$ tertutup dalam $Y$, sehingga
$\ker T=T^{-1}(\{\vc0\})$ tertutup dalam $X$. Karena $\vc0\in\ker T$,
\[
  N=\overline{\{\vc0\}}\subseteq\ker T.
\]
Maka rumus $\widetilde T([x])=T(x)$ terdefinisi dengan baik dan memberi
pemetaan linear unik $\widetilde T:X/N\to Y$ dengan
$T=\widetilde T\circ\pi$. Kontinuitas $\widetilde T$ mengikuti sifat
topologi hasil bagi: untuk setiap himpunan terbuka $O\subseteq Y$,
$\pi^{-1}(\widetilde T^{-1}(O))=T^{-1}(O)$ terbuka. Jadi $X/N$ bukan hanya
suatu ruang Hausdorff yang terkait, melainkan hasil bagi Hausdorff universal
dari~$X$.
\end{o001proof}
\end{o001solution}
